\documentclass[aspectratio=169]{beamer}
\input{header}
\coursecode{JKSSB}

\title{PYQ-Calibrated Final Retrieval \& Audit}
\subtitle{Lesson 58 --- Mixed Retrieval at Exam Depth}
\author{JKSSB Computer Awareness}
\date{\today}

\begin{document}
\frame{\titlepage}

\begin{frame}{Why This Capstone Exists}
  \begin{itemize}
    \item This is the \key{exam-eve} lesson: mixed retrieval across all units
      at the real JKSSB depth.
    \item Depth held across the drill set: \key{50\% L1} direct recall,
      \key{35\% L2} contrast and classification, \key{15\% L3} one-step
      application.
    \item Every item format JKSSB uses appears here: NOT/EXCEPT,
      near-neighbour, abbreviation, statement-evaluation, matching, and
      assertion--reason.
    \item Method first, facts second: \key{audit the stem}, then retrieve
      the fact.
  \end{itemize}
  \begin{block}{The one rule}
    Read the stem twice, name the format, then answer. Format errors cost
    more marks than fact errors.
  \end{block}
\end{frame}

\begin{frame}{The Six JKSSB Item Formats}
  \begin{center}
    \begin{tabular}{p{0.26\textwidth}p{0.62\textwidth}}
      \toprule
      \textbf{Format} & \textbf{What the stem looks like} \\
      \midrule
      NOT / EXCEPT & ``Which is NOT \dots'' --- three options true, one false \\
      Near-neighbour & Two confusable terms side by side (RAM/ROM, POP3/IMAP) \\
      Abbreviation & Expand the short form, or name the short form \\
      Statement-evaluation & Statements I/II/III; pick the correct combination \\
      Matching & Match List I with List II, or pick the correct pairing \\
      Assertion--reason & Assertion (A) + Reason (R); judge each, then the link \\
      \bottomrule
    \end{tabular}
  \end{center}
  \vspace{0.25cm}
  \muted{Each format has its own attack below. Never answer before naming it.}
\end{frame}

\begin{frame}{Format 1: NOT / EXCEPT Items}
  \begin{itemize}
    \item Circle the word \key{NOT}, \key{EXCEPT}, or \key{incorrect} before
      reading the options.
    \item Mark each option \key{True} or \key{False} in the margin; the odd
      one out is the answer.
    \item Trap: the three true options are all attractive. Do not pick the
      first true statement you see.
  \end{itemize}
  \begin{exampleblock}{Worked mini-example \textit{[Original]}, pattern after PYQ-03}
    Which is NOT true of cache memory? A. Holds recent data (true).
    B. Faster than RAM (true). C. Permanent backup of RAM (false).
    D. Exploits locality (true). Answer: \key{C} --- cache holds copies,
    not a backup.
  \end{exampleblock}
\end{frame}

\begin{frame}{Format 2: Near-Neighbour Distractors}
  \begin{itemize}
    \item For every key term, retrieve the triple: \key{what it is / what it
      is not / nearest confusion}.
    \item Classic pairs: RAM vs ROM, POP3 vs IMAP, Section Break vs Page
      Break, signature vs encryption, MICR vs OCR vs OMR.
    \item Read \key{both} words of the pair; the examiner swaps exactly one.
  \end{itemize}
  \begin{exampleblock}{Worked mini-example \textit{[Original]}, pattern after TRAP-15}
    ``A digital signature encrypts the mail body so only the recipient can
    read it.'' Wrong: signature gives \key{authenticity and integrity};
    \key{encryption} gives confidentiality. One verb swapped, whole item flips.
  \end{exampleblock}
\end{frame}

\begin{frame}{Format 3: Abbreviation and Shortcut Items}
  \begin{itemize}
    \item Give the \key{full form on first sight}: POP3 --- Post Office
      Protocol version 3; ISP --- Internet Service Provider.
    \item Shortcuts are exact: \key{F5} starts the slide show from the
      beginning; \key{F7} is spell check; a formula starts with \key{=}.
    \item Near-neighbours here are single-letter swaps: IMAP vs IMAC,
      POP3 vs SMTP. Spell the abbreviation out loud.
  \end{itemize}
  \begin{exampleblock}{Worked mini-example \textit{[Original]}, pattern after PYQ-34}
    Expand ISP. A. Internet Service Provider (correct). B. Intranet System
    Program. C. Internet Server Process. D. Internal Service Protocol.
    Only \key{A} matches the exact exam wording.
  \end{exampleblock}
\end{frame}

\begin{frame}{Format 4: Statement-Evaluation (I/II/III)}
  \begin{itemize}
    \item Judge \key{each statement alone} first: write T or F beside
      I, II, III.
    \item Then keep only the option whose combination matches your T/F
      pattern.
    \item Trap: one false statement hides inside two true ones (often the
      RAM-volatility or PC-address claim).
  \end{itemize}
  \begin{exampleblock}{Worked mini-example \textit{[Original]}, pattern after PYQ-03/PYQ-04}
    I. Registers are fastest (T). II. Cache is faster than RAM (T).
    III. RAM is non-volatile (F). Correct combination: \key{I and II only}.
  \end{exampleblock}
\end{frame}

\begin{frame}{Format 5: Matching Items}
  \begin{itemize}
    \item Match what you \key{know first}: lock one certain pair, strike
      out every option that breaks it.
    \item Device drills decide these: plotter is output, scanner is input,
      MICR reads magnetic ink, BIOS boots the machine.
    \item If two options survive, compare only the pairs that differ.
  \end{itemize}
  \begin{exampleblock}{Worked mini-example \textit{[Original]}, pattern after PYQ-08}
    Which pairing is correct? A. Plotter --- input (wrong). B. Scanner ---
    output (wrong). C. MICR --- magnetic ink characters (correct).
    D. Sound card --- input (wrong). Answer: \key{C}.
  \end{exampleblock}
\end{frame}

\begin{frame}{Format 6: Assertion--Reason Items}
  \begin{itemize}
    \item Step 1: judge (A) alone. Step 2: judge (R) alone. Step 3: ask
      whether (R) \key{explains} (A) --- not merely whether both are true.
    \item Both-true-but-unlinked is the standard trap; do not pick the
      ``explains'' option unless the link is direct.
  \end{itemize}
  \begin{exampleblock}{Worked mini-example \textit{[Original]}, pattern after PYQ-18}
    (A) IPv6 uses 128-bit addresses. (R) HTTP/3 runs over QUIC, not TCP.
    Both true, but (R) does not explain (A). Answer: \key{both true, (R) is
    not the explanation of (A)}.
  \end{exampleblock}
\end{frame}

\begin{frame}{The Negative-Stem Audit Method}
  \begin{itemize}
    \item Step 1: \key{underline} NOT, EXCEPT, incorrect, or CANNOT in the
      stem.
    \item Step 2: convert the stem to positive: ``find the three true
      statements, the leftover wins.''
    \item Step 3: test every option with a T/F mark; pick the single
      \key{False} one.
    \item Step 4: re-read the stem once more before marking the OMR ---
      confirm you answered the negative question.
  \end{itemize}
  \begin{alertblock}{Why this matters}
    With a \key{$-1/4$ penalty} per wrong answer, one misread NOT-stem
    costs the mark plus the penalty. Blank scores zero; a guess after
    eliminating two options is the only gamble worth taking.
  \end{alertblock}
\end{frame}

\begin{frame}{The Defective-Option Audit Method}
  \begin{itemize}
    \item Some official options are \key{defective}: loose wording the final
      key may not sustain. Quote them verbatim, flag them, and learn the
      stable concept instead.
    \item Defect 1, quoted verbatim: ``Program Counter stores the address of
      the \key{current} instruction being executed.'' \key{Flag:} the PC holds
      the address of the \key{next} instruction to fetch (TRAP-04).
    \item Defect 2, quoted verbatim (option wording): ``BCC hides sender and
      recipients from \key{each other}.'' \key{Flag:} BCC hides recipients
      from \key{each other}; the sender always knows, and To/Cc stay visible
      to all (TRAP-16).
  \end{itemize}
  \begin{block}{Stable concepts to carry}
    PC = next instruction. BCC = recipients hidden from each other only.
    The final key (Tier 1) always wins over any single option's wording.
  \end{block}
\end{frame}

\begin{frame}{Answer-Rationale Rules}
  \begin{itemize}
    \item Every answer needs two halves: \key{why the key is right} and
      \key{why each distractor is wrong}.
    \item Name the trap each distractor instantiates (e.g.\ ``reverses
      RAM/ROM,'' ``swaps Section/Page Break,'' ``states 256-bit IPv6'').
    \item Keep exam wording exact: full forms, shortcut keys, command
      spellings, and protocol names verbatim.
    \item One fact per line; never invent a paper number, page, or key date.
  \end{itemize}
  \begin{alertblock}{Rationale template}
    Correct option + one-line authority. Then, per distractor: what it
    claims, which trap it matches, and the one-line correction.
  \end{alertblock}
\end{frame}

\begin{frame}{Time and Penalty Discipline}
  \begin{itemize}
    \item Junior Assistant and Website Operator papers: \key{80 items, 80
      minutes, $-1/4$ per wrong answer}. Roughly one minute per item.
    \item First pass: answer every \key{L1 recall} item immediately; mark
      the rest.
    \item Second pass: L2 classification and matching with the T/F margin
      marks.
    \item Third pass: L3 scenarios only. \key{Blank} beats a blind guess;
      guess only after eliminating at least two options.
  \end{itemize}
\end{frame}

\begin{frame}{Quick-Hit Rail: Units 1--6}
  \begin{itemize}
    \item U1 Fundamentals: bit smallest; 10 = $1010_2$; ICs = third
      generation; BCD = 4 bits per digit.
    \item U2 CPU: PC = next instruction; interrupts at instruction
      boundaries; overflow is not carry-out; diagnose the stated bottleneck.
    \item U3 Memory: RAM volatile, ROM non-volatile; cache faster than RAM;
      CPU never touches secondary storage directly.
    \item U4 I/O: plotter/printer/sound card = output; scanner/light pen/
      digitizer = input; MICR = magnetic ink; Section Break changes
      headers/footers.
    \item U5 Software/OS: driver is software; OS is system software; BIOS
      boots; open-source = view/change/share.
    \item U6 Programming: compiler vs interpreter; stack is LIFO; algorithm
      before code.
  \end{itemize}
\end{frame}

\begin{frame}{Quick-Hit Rail: Units 7--11}
  \begin{itemize}
    \item U7 Office: formula starts with =; F5 from beginning, F7 spell
      check; Mail Merge = letters + database; Select query retrieves.
    \item U8 Internet/Email: @ separates user name and domain; POP3
      downloads, IMAP syncs; ISP = Internet Service Provider.
    \item U9 Networking: IPv6 = 128-bit; BGP is inter-domain; HTTP/3 uses
      QUIC over UDP; ICMP is control only.
    \item U10 Security: signature = authenticity, encryption =
      confidentiality; phishing vs ransomware vs DDoS kept distinct.
    \item U11 Governance: e-NAM = agricultural marketing; G2C/G2G/G2B/G2E
      expansions; Digital India and NeGP/e-Kranti roles.
  \end{itemize}
\end{frame}

\begin{frame}{Lesson 58: Recall}
  \begin{itemize}
    \item Name the format first: NOT/EXCEPT, near-neighbour, abbreviation,
      statement-evaluation, matching, assertion--reason.
    \item Negative stems: underline NOT, mark T/F per option, re-read before
      marking.
    \item Defective options: quote verbatim, flag, and carry the stable
      concept (PC = next; BCC = recipients hidden from each other).
    \item Every rationale: why the key is right plus why each distractor is
      wrong, with the trap named.
    \item Hold the 50/35/15 depth; spend time where L2/L3 traps cluster, and
      respect the $-1/4$ penalty.
  \end{itemize}
\end{frame}
\end{document}
